
%\subsection{Collective Epistemic Dynamics)}

\subsubsection{ Purpose and Scope}

This appendix extends  Epistemic dynamics to \emph{collective} epistemic
systems. Rather than a single agent executing a single FRFP pipeline, we
consider populations of agents who:

\begin{itemize}
  \item share an explicit artifact space $E$,
  \item each carry their own tacit state and degradation dynamics,
  \item exchange only explicit artifacts (no tacit information flow),
  \item may be coupled by explicit consensus mechanisms and institutions.
\end{itemize}

The objective is to characterize:

\begin{enumerate}
  \item how to view a population as a single large FRFP system,
  \item how stochastic collective evolution can be modeled as a Markov process
        on this system,
  \item how hallucination and other failures lift from individuals to
        populations, and
  \item which collective design changes \emph{cannot}, even in principle,
        eliminate hallucination.
\end{enumerate}

This section provides the bridge between individual pipelines (Sections A.2–A.8),
impossibility results (Appendix A.11), and institutional design (Appendix B and Appendix C).

\subsubsection{Epistemic Agents}

\begin{definition}[Epistemic agent]
An \emph{epistemic agent} is a tuple
\[
  A_i = (T_i, \,\preceq_i, \,\delta_i, \,T^e_i),
\]
where:
\begin{itemize}
  \item $T_i$ is a tacit context space equipped with a preorder
        $\preceq_i$ encoding ``at least as grounded as'';
  \item $\delta_i : T_i \to T_i$ is a (tacit) context degradation operator;
  \item $T^e_i := T_i \times \mathbb{R}_{\ge 0}$ is the extended tacit state
        space, pairing grounding $t_i \in T_i$ with tacit credence
        $c_i \in \mathbb{R}_{\ge 0}$.
\end{itemize}
\end{definition}

Intuitively, $A_i$ captures a single human (or human-like) reasoner as modeled
in Epistemic Dynamics: grounding lives in $T_i$, weakens under $\delta_i$,
and credence lives in the additional component of $T^e_i$.

\subsubsection{Epistemic Populations}

\begin{definition}[Epistemic population and category \textbf{Pop}]
An \emph{epistemic population} is a finite or countable family
$P = (A_i)_{i \in I}$ of epistemic agents indexed by a set $I$.
The category $\mathbf{Pop}$ of epistemic populations is defined by:
\begin{itemize}
  \item objects: epistemic populations $P = (A_i)_{i \in I}$;
  \item morphisms $F : P \to P'$: pairs $(f,(F_i)_{i\in I})$ where
        $f : I \to I'$ is a function on agent indices and, for each $i$,
        \[
          F_i : T_i \to T'_{f(i)}
        \]
        is monotone with respect to $\preceq_i,\preceq'_{f(i)}$ and
        commutes with degradation and credence:
        \[
          F_i \circ \delta_i = \delta'_{f(i)} \circ F_i.
        \]
\end{itemize}
\end{definition}

Morphisms in $\mathbf{Pop}$ model re-indexing or restructuring of populations:
embedding agents into larger organizations, replacing agents, or transferring
tacit states along institutional lines.

\subsubsection{Shared Explicit Reasoning Category}

Let $\mathcal{C} := N \ltimes E$ be the reasoning category of Appendix~A,
where $E$ denotes explicit artifacts and $N$ navigation moves in explicit
space. For a fixed index set $I$, we form the product category
\[
  \mathcal{C}^I,
\]
whose objects are $I$-indexed families $X = (X_i)_{i \in I}$ of explicit
states and whose morphisms are families of explicit updates. Intuitively,
$\mathcal{C}^I$ is the ``big explicit state space'' of the whole population:
it records, for each agent, which task and explicit artifacts they are
currently working with.

We assume a single shared explicit artifact space $E$; each agent $i$ is
equipped with a representational backflow map
\[
  \mathrm{RB}_i : E \to T_i,
\]
as in Appendix~A. Explicit artifacts are globally visible; tacit grounding
remains agent-local.

\subsubsection{Communication Structure}

\begin{definition}[Communication graph]
A \emph{communication structure} on population $P = (A_i)_{i \in I}$ is a
directed graph
\[
  G = (I,\;\to),
\]
where $i \to j$ indicates that explicit artifacts produced by $i$ may be
consumed by $j$ (for example, by reading a paper, listening to a briefing, or
reviewing a log). No tacit information is transmitted across $G$:
only explicit artifacts traverse edges, and representational backflow
$\mathrm{RB}_j$ determines how they affect $T_j$.
\end{definition}

\subsubsection{Collective Dynamics as a Markov System}

A \emph{collective execution} in earlier drafts was modeled as a bare
sequence of events $(i_n,e_n) \in I \times E$ indicating ``agent $i_n$
produces artifact $e_n$''. We now package such executions into a stochastic
dynamical system over the large explicit category $\mathcal{C}^I$.

\begin{definition}[Collective state]
A collective explicit state is an object $X \in \mathcal{C}^I$ together with
a population $P = (A_i)_{i\in I}$ and communication graph $G=(I,\to)$.
We write $(P,G;X)$ for such a state.
\end{definition}

\begin{definition}[Population semantics]
A \emph{population semantics functor} is a functor
\[
  \overline{J} : \mathcal{C}^I \to \prod_{i\in I} T^e_i
\]
assigning to each collective explicit state $X$ the tuple of extended tacit
states $\bigl((t_i,c_i)\bigr)_{i\in I}$ obtained when agents in $P$ consume
the artifacts in $X$ (subject to $G$) and update using the operators of
Appendix~B (degradation, confidence inflation, requirement escalation, and
feedback-coupled degradation).
\end{definition}

\begin{definition}[Markovian collective dynamics]
A \emph{Markovian collective dynamics} on $(P,G)$ is a transition kernel
\[
  M : \mathrm{Obj}(\mathcal{C}^I) \to \mathcal{P}\bigl(\mathrm{Obj}(\mathcal{C}^I)\bigr)
\]
assigning to each collective state $X$ a probability measure $M(X,\cdot)$
on next states $Y$, such that each transition can be factored as:
\begin{enumerate}
  \item select an acting agent $i_n \in I$ and explicit move in $\mathcal{C}$;
  \item update $X$ to $Y$ accordingly (producing a new artifact $e_n$ in $E$);
  \item for each $j$ with $i_n \to j$, update $(t_j,c_j)$ via the
        feedback-coupled dynamics of Appendix~B using $e_n$;
  \item apply baseline degradation for all agents.
\end{enumerate}
Equivalently, one may view $M$ as a functor
\[
  M : \mathcal{C}^I \longrightarrow \mathbf{Kl}(D),
\]
where $\mathbf{Kl}(D)$ is the Kleisli category of the discrete distribution
monad $D$, but we will only use the induced transition kernel on objects.
\end{definition}

A \emph{collective execution} is then a random path
\[
  X_0, X_1, X_2, \dots
\]
in $\mathcal{C}^I$ sampled from $M$ starting at some initial state $X_0$,
with tacit trajectories $\overline{J}(X_n) = \bigl((t_{i,n},c_{i,n})\bigr)_{i\in I}$.

\subsubsection{Collective Hallucination and Stopping Times}

We now lift the notion of hallucination to the population level.

\begin{definition}[Collective hallucination]
Given a population $P=(A_i)_{i\in I}$, a communication graph $G$, and a
collective execution $(X_n)_{n\ge 0}$ with associated tacit states
$(t_{i,n},c_{i,n})$, a \emph{collective hallucination} occurs at step $n$
if there exists an explicit artifact $e \in E$ and a nontrivial subset
$J \subseteq I$ such that:
\begin{itemize}
  \item each $j \in J$ has accepted or propagated $e$ (in the sense of
        Appendix~B: $e$ enters their explicit pipeline), and
  \item for each $j \in J$,
        \[
          t_{j,n} \prec_j \mathrm{Req}_j(e,c_{j,n}),
        \]
        i.e., the available tacit grounding is strictly below the requirement
        demanded by $e$ in context $c_{j,n}$.
\end{itemize}
\end{definition}

Let $\Omega$ be the space of all collective paths $(X_n)_{n\ge 0}$ induced
by $M$ from $X_0$, and let $P_M$ be the corresponding path measure.

\begin{definition}[Collective hallucination time]
The \emph{collective hallucination time} $\tau^{\mathrm{coll}} :
\Omega \to \mathbb{N} \cup \{\infty\}$ is the stopping time
\[
  \tau^{\mathrm{coll}}(\omega)
  := \inf\{\,n \ge 0 : \omega \text{ exhibits a collective hallucination at step } n\,\},
\]
with the convention that $\inf \varnothing = \infty$.
\end{definition}

Under the Epistemic Dynamics assumptions (finite tacit capacity, non-invertible
degradation, confidence inflation, requirement escalation, and
feedback-coupled degradation), individual agents exhibit almost-sure
hallucination along sufficiently long admissible runs (Appendix~B).
The purpose of this appendix is to show how such individual inevitability
lifts to $\tau^{\mathrm{coll}}$ at the population level.

\subsubsection{Consensus as an Explicit Functor}

Consensus mechanisms act purely in explicit space.

\begin{definition}[Consensus functor]
A \emph{consensus mechanism} on the explicit category $\mathcal{C}$ is a
functor
\[
  C : \mathcal{C}^n \to \mathcal{C}
\]
for some finite arity $n \ge 2$, or more simply
\[
  C : E^n \to E
\]
on the underlying artifact space, modeling averaging, voting, or more
complicated aggregation pipelines. Consensus functors operate entirely in
explicit space: they take explicit artifacts as inputs and produce explicit
artifacts as outputs, without access to tacit states or requirements.
\end{definition}

\subsubsection{Functorial Consensus Layer}

Given a population $P=(A_i)_{i\in I}$ and communication structure $G$,
a consensus mechanism $C$ induces an explicit functor
\[
  C_P : E^n \to E
\]
that may be interleaved with ordinary explicit computation in
$\mathcal{C}^I$. We refer to such composites as the \emph{consensus layer}.
No consensus functor has access to the population semantics
$\overline{J}$ or to individual requirement functions $\mathrm{Req}_i$.

\subsubsection{Tacit-Stability of Consensus}

\begin{definition}[Tacit-stable consensus]
A consensus functor $C$ is \emph{tacit-stable} if, for every agent $i$ and
any inputs $e_1,\dots,e_n \in E$, whenever $e_1,\dots,e_n$ are all
individually below agent $i$'s requirement,
\[
  t_i \prec_i \mathrm{Req}_i(e_k,c_i)
  \quad \forall k \in \{1,\dots,n\},
\]
then the aggregate output $C(e_1,\dots,e_n)$ is also below requirement:
\[
  t_i \prec_i \mathrm{Req}_i\bigl(C(e_1,\dots,e_n),c_i\bigr).
\]
\end{definition}

\begin{theorem}[Consensus stability]
Every consensus functor that acts purely in explicit space is tacit-stable.
\end{theorem}

\begin{proof}[Proof sketch]
Consensus functors are compositions of explicit morphisms in $\mathcal{C}$.
By Explicit--Tacit Separation (ETS), explicit morphisms cannot increase
tacit grounding: they do not alter $t_i$, only what artifacts are presented
to representational backflow. If an agent's tacit state $t_i$ is below the
requirement for each input artifact, it remains below the requirement for
any artifact obtained by explicit-only transformation of those inputs. Thus
consensus cannot manufacture sufficient grounding from insufficient inputs.
\end{proof}

Tacit-stability says that explicit consensus can neither repair nor upgrade
tacit adequacy: it preserves ``being under-grounded'' whenever that is the
case for all inputs.

\subsubsection{Limits of Automated Consensus}

\begin{theorem}[No fully automated consensus elimination]
No fully automated consensus process acting purely in explicit space can
eliminate collective hallucination across all admissible populations and
collective dynamics.
\end{theorem}

\begin{proof}[Proof sketch]
By tacit-stability, consensus cannot transform artifacts that are
under-grounded for all relevant agents into artifacts that are adequately
grounded. Under the Epistemic Dynamics degradation assumptions, individual hallucination
is inevitable along sufficiently long runs. Since consensus mechanisms have no
access to tacit state, they cannot detect or remove these failures at the
collective level; at best they can smooth, average, or denoise explicit
behavior. Any claim of global hallucination-elimination by explicit-only
consensus would contradict the Epistemic Dynamics inevitability results when applied
agent-wise.
\end{proof}

Automation may reduce variance or noise in explicit artifacts, but it cannot
introduce new tacit grounding or globally halt its degradation.

\subsubsection{Collective Inevitability}

\begin{theorem}[Collective hallucination inevitability]
Consider a population $P=(A_i)_{i\in I}$ of epistemic agents with:
\begin{itemize}
  \item finite or countable index set $I$,
  \item finite tacit capacity for each $A_i$,
  \item non-invertible degradation $\delta_i$,
  \item confidence inflation and requirement escalation as in Appendix~B.
\end{itemize}
Let $M$ be any Markovian collective dynamics on $\mathcal{C}^I$ that
respects these per-agent dynamics and produces infinitely many
surface-plausible explicit artifacts along almost every path.
Then the collective hallucination time $\tau^{\mathrm{coll}}$ is almost
surely finite:
\[
  P_M\bigl(\tau^{\mathrm{coll}} < \infty\bigr) = 1.
\]
Moreover, this property is invariant under the insertion of explicit-only
consensus functors: composing $M$ with any finite consensus layer that acts
solely in explicit space does not change the almost-sure finiteness of
$\tau^{\mathrm{coll}}$.
\end{theorem}

\begin{proof}[Proof sketch]
By Appendix~B, each agent $A_i$ exhibits almost-sure hallucination along
sufficiently long admissible runs of any FRFP pipeline that continues to
produce surface-plausible artifacts. The collective Markov dynamics $M$
induces, for each $i$, a marginal stochastic process on $T^e_i$ with the
same qualitative properties. Since the population is finite or countable,
and collective execution produces infinitely many plausible artifacts with
nonzero probability of being consumed by each agent (via $G$), the event
that at least one agent hallucinates in finite time has probability~$1$.
The definition of collective hallucination requires only that a nontrivial
subset of agents jointly accept an under-grounded artifact, which occurs
almost surely given enough time and interaction. Tacit-stability of
consensus implies that explicit aggregation cannot remove such events once
they arise; it can at best transform one hallucinated artifact into another.
\end{proof}

\subsubsection{Structural Implications of Collective Dynamics}

The population-level formalism above does more than restate
inevitability in a collective setting. It yields several structural notions
that are used in later sections.

\paragraph{(1) Equivalence classes of collective designs.}
Viewing populations as objects in $\mathbf{Pop}$ and their joint explicit
reasoning as trajectories in $\mathcal{C}^I$, we can treat different UIs,
workflows, and consensus wrappers as endofunctors on $\mathcal{C}^I$ that
preserve the population semantics $\overline{J}$. Designs that differ only
by such explicit refactorings are \emph{FRFP-equivalent}: they have the same
tacit semantics and the same hallucination-time distribution, even if their
surface appearance differs. This provides a criterion for when two collective
protocols are ``the same thing in principle''.

\paragraph{(2) Refinement order on collective systems.}
Given two Markovian collective dynamics $M_1,M_2$ on the same population,
with corresponding collective hallucination times
$\tau^{\mathrm{coll}}_1,\tau^{\mathrm{coll}}_2$, we may compare their
finite-horizon risks
\[
  p_N^{(k)} := P_{M_k}(\tau^{\mathrm{coll}}_k \le N), \quad k\in\{1,2\}.
\]
We say that $M_2$ is a \emph{dynamic refinement} of $M_1$ if
$p_N^{(2)} \le p_N^{(1)}$ for all $N$. This yields a partial order on
collective implementations: some architectures are strictly safer than
others, but purely explicit-only transformations that preserve semantics
cannot improve this order.

\paragraph{(3) Compositional risk analysis.}
Because the collective dynamics lives in $\mathcal{C}^I$ and is built from
per-agent Epistemic Dynamics, we can study how subsystem risks compose.
For example, a scientific subsystem, an engineering subsystem, and a policy
subsystem can each be modeled as a subpopulation with its own hazard profile;
the global population then inherits a collective hallucination time equal to
the minimum of the subsystem times. This provides a formal basis for
``defense in depth'' arguments.

\paragraph{(4) Limits of monitoring and metrics.}
Collective executions produce rich explicit traces: logs, dashboards,
metrics, and aggregated indicators. In this formalism, these traces live in
$\mathcal{C}^I$ or in functorial images thereof. The population semantics
$\overline{J}$ and individual requirements $\mathrm{Req}_i$, however, remain
tacit. Many distinct tacit configurations can induce the same distribution
over explicit traces. Hence correctness and hallucination risk are
\emph{non-identifiable from logs alone}: explicit-only monitoring cannot, in
general, determine how well grounded a population's beliefs really are.

\paragraph{(5) Scaling and phase transitions.}
By parametrizing collective systems by population size $|I|$, communication
structure $G$, and per-agent degradation properties, Appendix~A.10 provides a
language for asking how epistemic risk scales with size and complexity.
Nothing in the formalism guarantees that larger or more connected populations
are safer; on the contrary, the inevitability theorem applies uniformly once
the Epistemic Dynamics assumptions hold. This sets the stage for Appendix~A.12, which
studies which institutional forms can slow or bound such risks across
populations and time.

Taken together, these structures show that Collective Epistemic Dynamics is not merely a
corollary of individual hallucination: it is a genuinely new level of
description in which equivalence, refinement, observability, and scaling of
\emph{collective} epistemic systems can be stated and analyzed formally.
